%%% FIELD proof-prop-endpoint-reductions
Throughout, the potential-outcome schedule is fixed and the only expectation is over the assignment law, as required by \(\texttt{ass:design-exogeneity}\) and \(\texttt{def:randomization-risk}\).

First let \(0\le s\le s'\le J\).  If \(Y\in\Theta_{J,s}\), then \(D_a(Y)\le s\le s'\) for each \(a\in\{0,1\}\).  Hence, directly from \(\texttt{def:bounded-mismatch-class}\),
\[
 \Theta_{J,s}\subseteq\Theta_{J,s'}.
\]
The condition in \(\texttt{def:design-class}\) restricts a law only to be selected ex ante and independently of \(Y\); the numerical sensitivity label may inform which law is selected but does not remove any outcome-independent assignment law from the feasible simplex.  Thus \(\mathfrak Q_{J,s}=\mathfrak Q_{J,s'}\), while \(\mathfrak D_J\) is the same set at both sensitivity values by \(\texttt{def:estimator-class}\).  Consequently, for every feasible \((Q,\delta)\),
\[
 \sup_{Y\in\Theta_{J,s}}\mathcal R_{J,s}(Q,\delta;Y)
 \le
 \sup_{Y\in\Theta_{J,s'}}\mathcal R_{J,s'}(Q,\delta;Y),
\]
where the two risk formulas agree pointwise by \(\texttt{def:randomization-risk}\).  Taking the infimum over the common feasible procedure class and applying \(\texttt{def:minimax-risk}\) gives
\[
 R^*(J,s)\le R^*(J,s').
\]

Now set \(s=0\).  By \(\texttt{def:bounded-mismatch-class}\), every \(Y\in\Theta_{J,0}\) has constants \(y_j(0),y_j(1)\in\{0,1\}\) such that
\[
 Y_{j1}(a)=Y_{j2}(a)=y_j(a),
 \qquad a\in\{0,1\},\quad 1\le j\le J.
\]
Under \(q^{\mathrm{pair}}\), exactly one member of each pair is treated.  Define the measurable estimator
\[
 \widehat\tau_{\mathrm{pair}}(A,Y^{\mathrm{obs}})
 =\frac1J\sum_{j=1}^J\sum_{r=1}^2
 \bigl\{A_{jr}Y^{\mathrm{obs}}_{jr}-(1-A_{jr})Y^{\mathrm{obs}}_{jr}\bigr\}.
\]
Each pair contributes a number in \([-1,1]\), so \(\widehat\tau_{\mathrm{pair}}\in[-1,1]\) and therefore belongs to \(\mathfrak D_J\).  Pair homogeneity and consistency in the definition of \(Y^{\mathrm{obs}}\) imply, assignment by assignment,
\[
 \widehat\tau_{\mathrm{pair}}
 =\frac1J\sum_{j=1}^J\{y_j(1)-y_j(0)\}
 =\frac1{2J}\sum_{j=1}^J\sum_{r=1}^2
 \{Y_{jr}(1)-Y_{jr}(0)\}
 =\tau(Y),
\]
the last equality being \(\texttt{def:sate}\).  Thus the worst-case randomization risk of this procedure is zero.  Squared loss is nonnegative, so \(\texttt{def:minimax-risk}\) yields
\[
 R^*(J,0)=0.
\]

It remains to show that the optimal invariant orbit law is unique at this endpoint.  Let \(q\in\Delta(\mathcal N_J)\) place positive mass on some \(n\ne(0,0,J)\).  Then every assignment in orbit \(n\) contains an unmixed pair.  Choose one labelled assignment \(a\) in that orbit and one such pair \(j\).  By the orbit construction, \(Q(a)=q_n/C_n>0\).  Let \(Y^{(0)}\) be the all-zero schedule.  If pair \(j\) is all control under \(a\), obtain \(Y^{(1)}\) by setting \(Y^{(1)}_{j1}(1)=Y^{(1)}_{j2}(1)=1\) and leaving all other potential outcomes zero.  If pair \(j\) is all treated, instead set \(Y^{(1)}_{j1}(0)=Y^{(1)}_{j2}(0)=1\).  In either case, both schedules lie in \(\Theta_{J,0}\), and they produce the same observed data under \(a\).  Their SATEs are respectively zero and \(\varepsilon/J\), where \(\varepsilon=1\) in the all-control case and \(\varepsilon=-1\) in the all-treated case.  Therefore, if \(d\) is the common estimator value at that observed-data cell,
\[
 \max\left\{d^2,\left(d-\frac{\varepsilon}{J}\right)^2\right\}
 \ge \frac1{4J^2}.
\]
For at least one of \(Y^{(0)}\) and \(Y^{(1)}\), nonnegativity of all the other assignment contributions then gives
\[
 \mathcal R_{J,0}(Q,\delta;Y)
 \ge \frac{q_n}{4C_nJ^2}>0.
\]
Thus no such \(q\) can attain the already established value zero.  The only remaining orbit law is \(\delta_{(0,0,J)}=q^{\mathrm{pair}}\), so \(\texttt{def:finite-optimal-design-set}\) gives
\[
 \mathcal Q^*(J,0)=\{q^{\mathrm{pair}}\}.
\]

Finally set \(s=J\).  Every binary schedule has \(D_a(Y)\le J\), and hence
\[
 \Theta_{J,J}=\{0,1\}^{\mathcal P_J\times\{0,1\}}
\]
by \(\texttt{def:bounded-mismatch-class}\).  This is exactly the unrestricted binary fixed-population class studied by \(\cite{SudijonoDobribanTchetgen2026}\).  Since the design, estimator, target, and loss in \(\texttt{def:unrestricted-endpoint-risk}\) are the same as those in \(\texttt{def:minimax-risk}\) at \(s=J\), their defining infimum--supremum problems coincide, and therefore
\[
 R^*(J,J)=r^{\mathrm U}_{2J}.
\]

For clarity, this binary-class identity is distinct from the comparison with \(\cite{Hull2026BoundedOutcomes}\).  Let
\[
 \mathcal Y_{\mathrm{bin}}=\{0,1\}^{\mathcal P_J\times\{0,1\}},
 \qquad
 \mathcal Y_{[0,1]}=[0,1]^{\mathcal P_J\times\{0,1\}},
\]
and let \(V_{[0,1]}\) denote Hull's minimax value on \(\mathcal Y_{[0,1]}\).  Then \(\mathcal Y_{\mathrm{bin}}\subsetneq\mathcal Y_{[0,1]}\): the former is only the endpoint-valued subclass of the latter.  By \(\texttt{lem:unrestricted-binary-bernoulli-minimax}\), Hull's independent \(\operatorname{Bernoulli}(1/2)\) design and nonlinear posterior-mean rule, denoted \((Q_H,\delta_H)\), obey the uniform bounded-interval upper bound
\[
 \sup_{Y\in\mathcal Y_{[0,1]}}
 \mathcal R(Q_H,\delta_H;Y)=V_{[0,1]}.
\]
Restriction of this one procedure to binary observations is feasible for the binary problem, so
\[
 r^{\mathrm U}_{2J}
 \le \sup_{Y\in\mathcal Y_{\mathrm{bin}}}
 \mathcal R(Q_H,\delta_H;Y)
 \le V_{[0,1]}.
\]
The same cited lemma supplies a least-favourable prior \(\Pi_H\) supported on endpoint schedules, hence supported on \(\mathcal Y_{\mathrm{bin}}\), whose optimal Bayes risk is \(V_{[0,1]}\).  For every binary procedure \((Q,\delta)\), the supremum dominates its \(\Pi_H\)-average risk.  Taking infima gives the reverse inequality
\[
 r^{\mathrm U}_{2J}
 =\inf_{Q,\delta}\sup_{Y\in\mathcal Y_{\mathrm{bin}}}
 \mathcal R(Q,\delta;Y)
 \ge
 \inf_{Q,\delta}\int_{\mathcal Y_{\mathrm{bin}}}
 \mathcal R(Q,\delta;Y)\,d\Pi_H(Y)
 =V_{[0,1]}.
\]
Thus the two values are equal and \((Q_H,\delta_H)\), restricted to binary data, attains \(r^{\mathrm U}_{2J}\).  This is a uniform-upper-bound plus endpoint-prior-lower-bound sandwich; it does not identify the binary and bounded-interval parameter classes or their experiments.

Under independent \(\operatorname{Bernoulli}(1/2)\) assignment, every labelled assignment has probability \(4^{-J}\).  Orbit \(n\) contains \(C_n\) assignments, so its mass is
\[
 Q_H\{A\text{ lies in orbit }n\}=\frac{C_n}{4^J}=q_n^{\mathrm{Ber}}
\]
by \(\texttt{def:canonical-segment}\).  The binary schedule class and squared-error loss are invariant under \(G_J\); averaging the optimal procedure over this finite group leaves the already invariant Bernoulli design unchanged and, by convexity of squared loss, cannot increase its maximum risk.  The lossless invariant reduction in \(\texttt{thm:orbit-game-face}\) therefore supplies an invariant orbit estimator \(\eta\) achieving \(R^*(J,J)\) with orbit law \(q^{\mathrm{Ber}}\).  By \(\texttt{def:finite-optimal-design-set}\),
\[
 q^{\mathrm{Ber}}\in\mathcal Q^*(J,J),
\]
and the unrestricted attaining procedure is the asserted nonlinear posterior-mean rule.
